%%%%%%%% ICML 2026 EXAMPLE LATEX SUBMISSION FILE %%%%%%%%%%%%%%%%%

\documentclass{article}

% Recommended, but optional, packages for figures and better typesetting:
\usepackage{microtype}
\usepackage{graphicx}
\usepackage{subcaption}
\usepackage{booktabs} % for professional tables

% hyperref makes hyperlinks in the resulting PDF.
% If your build breaks (sometimes temporarily if a hyperlink spans a page)
% please comment out the following usepackage line and replace
% \usepackage{icml2026} with \usepackage[nohyperref]{icml2026} above.
\usepackage{hyperref}


% Attempt to make hyperref and algorithmic work together better:
\newcommand{\theHalgorithm}{\arabic{algorithm}}

% Use the following line for the initial blind version submitted for review:
\usepackage{icml2026}

% For preprint, use
% \usepackage[preprint]{icml2026}

% If accepted, instead use the following line for the camera-ready submission:
% \usepackage[accepted]{icml2026}

\usepackage{amsmath}
\usepackage{amssymb}
\usepackage{mathtools}
\usepackage{amsthm}
\usepackage{booktabs,multirow}
\usepackage{tikz}
\usetikzlibrary{shapes,arrows,positioning,calc}
\usepackage{mdframed}
\usepackage{listings}
\lstset{
  basicstyle=\small\ttfamily,
  breaklines=true,
  breakatwhitespace=true,
  columns=flexible,
  keepspaces=true
}

% if you use cleveref..
\usepackage[capitalize,noabbrev]{cleveref}

%%%%%%%%%%%%%%%%%%%%%%%%%%%%%%%%
% THEOREMS
%%%%%%%%%%%%%%%%%%%%%%%%%%%%%%%%
\theoremstyle{plain}
\newtheorem{theorem}{Theorem}[section]
\newtheorem{proposition}[theorem]{Proposition}
\newtheorem{lemma}[theorem]{Lemma}
\newtheorem{corollary}[theorem]{Corollary}
\theoremstyle{definition}
\newtheorem{definition}[theorem]{Definition}
\newtheorem{assumption}[theorem]{Assumption}
\theoremstyle{remark}
\newtheorem{remark}[theorem]{Remark}

% Todonotes is useful during development; simply uncomment the next line
%    and comment out the line below the next line to turn off comments
%\usepackage[disable,textsize=tiny]{todonotes}
\usepackage[textsize=tiny]{todonotes}

% The \icmltitle you define below is probably too long as a header.
% Therefore, a short form for the running title is supplied here:
\icmltitlerunning{Submission and Formatting Instructions for ICML 2026}

\begin{document}

\twocolumn[
  \icmltitle{Expert-Guided Agents for Kernel Generation}

  % It is OKAY to include author information, even for blind submissions: the
  % style file will automatically remove it for you unless you've provided
  % the [accepted] option to the icml2026 package.

  % List of affiliations: The first argument should be a (short) identifier you
  % will use later to specify author affiliations Academic affiliations
  % should list Department, University, City, Region, Country Industry
  % affiliations should list Company, City, Region, Country

  % You can specify symbols, otherwise they are numbered in order. Ideally, you
  % should not use this facility. Affiliations will be numbered in order of
  % appearance and this is the preferred way.
  \icmlsetsymbol{equal}{*}

  \begin{icmlauthorlist}
    \icmlauthor{Firstname1 Lastname1}{equal,yyy}
    \icmlauthor{Firstname2 Lastname2}{equal,yyy,comp}
    \icmlauthor{Firstname3 Lastname3}{comp}
    \icmlauthor{Firstname4 Lastname4}{sch}
    \icmlauthor{Firstname5 Lastname5}{yyy}
    \icmlauthor{Firstname6 Lastname6}{sch,yyy,comp}
    \icmlauthor{Firstname7 Lastname7}{comp}
    %\icmlauthor{}{sch}
    \icmlauthor{Firstname8 Lastname8}{sch}
    \icmlauthor{Firstname8 Lastname8}{yyy,comp}
    %\icmlauthor{}{sch}
    %\icmlauthor{}{sch}
  \end{icmlauthorlist}

  \icmlaffiliation{yyy}{Department of XXX, University of YYY, Location, Country}
  \icmlaffiliation{comp}{Company Name, Location, Country}
  \icmlaffiliation{sch}{School of ZZZ, Institute of WWW, Location, Country}

  \icmlcorrespondingauthor{Firstname1 Lastname1}{first1.last1@xxx.edu}
  \icmlcorrespondingauthor{Firstname2 Lastname2}{first2.last2@www.uk}

  % You may provide any keywords that you find helpful for describing your
  % paper; these are used to populate the "keywords" metadata in the PDF but
  % will not be shown in the document
  \icmlkeywords{Machine Learning, ICML}

  \vskip 0.3in
]

% this must go after the closing bracket ] following \twocolumn[ ...

% This command actually creates the footnote in the first column listing the
% affiliations and the copyright notice. The command takes one argument, which
% is text to display at the start of the footnote. The \icmlEqualContribution
% command is standard text for equal contribution. Remove it (just {}) if you
% do not need this facility.

% Use ONE of the following lines. DO NOT remove the command.
% If you have no special notice, KEEP empty braces:
\printAffiliationsAndNotice{}  % no special notice (required even if empty)
% Or, if applicable, use the standard equal contribution text:
% \printAffiliationsAndNotice{\icmlEqualContribution}

\begin{abstract}
High-performance GPU kernels are critical for reducing the exponentially growing computational costs of large language models, but their development heavily relies on manual tuning by domain experts. Although recent advances in LLMs show promise for automating kernel generation, existing methods struggle with both correctness and performance. They lack domain-specific optimization knowledge and rely on coarse-grained performance feedback, making it challenging to navigate the vast optimization space. 
We propose EGG, an \underline{E}xpert-\underline{G}uided Agent Framework for Kernel \underline{G}eneration. By integrating expert optimization knowledge, we decompose kernel generation into two hierarchical stages: algorithmic structure design establishes a high-quality foundation, followed by hardware-specific tuning that performs targeted optimizations through parallel mapping, tensor tiling, and memory optimization. This decomposition provides explicit optimization objectives at each stage, guiding the LLM to achieve stable and cumulative improvements. 
To support this workflow, a stage-aware multi-agent system coordinates specialized agents within stages and selectively propagates decisions across stages.
Evaluated on KernelBench and real-world workloads, our approach achieves 100\% correctness with 2.13$\times$ average speedup over PyTorch, outperforming existing LLM-based methods and compiler framework.
\end{abstract}

\section{Instruction}
The computational demands of large language models (LLMs) are growing exponentially, driving training and inference costs to millions of dollars. High-performance GPU kernels are critical to reducing these costs. However, kernel development heavily relies on manual tuning by domain experts. As GPU architectures evolve rapidly and diversify, this manual optimization process has become increasingly complex and expensive, making automated GPU kernel generation a critical research direction for large-scale machine learning.

Recent advances in LLMs have demonstrated strong performance on general-purpose code generation tasks (e.g., Python, C++), making them promising tools for GPU kernel generation. Existing approaches can be categorized into two classes: (1) specializing LLMs through fine-tuning or reinforcement learning (RL), and (2) building agent systems based on general-purpose LLMs.

Although RL has achieved success in decision-making tasks, its applicability to GPU kernel generation remains limited. High-quality kernel data are scarce, and kernel code must satisfy strict syntactic and semantic constraints. As a result, generated kernels often fail to compile or produce incorrect results. For example, AutoTriton reports an average correctness rate of less than 40\%.

% In contrast, agent approaches built on general-purpose LLMs avoid expensive training cost and can directly leverage pretrained models' strong code generation capabilities. Through multi-turn iterative refinement, these methods significantly improve correctness, offering better engineering feasibility and scalability. However, most existing agent systems rely on coarse-grained performance feedback to improve kernels and lack domain-specific kernel optimization knowledge. Without explicit guidance on what to optimize and how, agents often make random or conflicting modifications across iterations, resulting in limited performance gains. For example, GEAK shows little performance improvement after four optimization rounds.
In contrast, agent approaches built on general-purpose LLMs avoid expensive training costs and directly leverage pretrained models' strong code generation capabilities. Through multi-turn iterative refinement, these methods significantly improve correctness, offering better engineering feasibility and scalability. However, most existing agent systems lack domain-specific kernel optimization knowledge and rely on coarse-grained performance feedback (e.g., execution time). This causes agents to make blind trial-and-error modifications, resulting in limited performance gains. For example, GEAK shows little improvement after four optimization rounds.

The root cause lies in the intrinsic complexity of the kernel optimization space. Kernel performance is determined by intricate interactions among thread scheduling, memory hierarchy utilization, synchronization mechanisms, and hardware-specific features. Even minor code modifications can lead to orders-of-magnitude performance differences. Previous work has shown that even a small neural network subgraph on GPUs can exhibit an optimization space as large as $10^9$ configurations (Zhai et al., 2024). Without expert knowledge to systematically navigate this vast landscape, agents struggle to discover high-performance solutions.

Based on these observations, we argue that fully unlocking the potential of LLMs for kernel generation requires integrating expert-level kernel optimization knowledge into agent systems to guide model decisions. To this end, we propose XX, an expert-guided agent framework that automatically generates high performance GPU kernels.

Specifically, we introduce expert-guided staged optimization, decomposing the kernel generation workflow into two hierarchical levels: high-level algorithmic structure design and fine-grained hardware-specific tuning. In the algorithmic structure design stage, we combine multi-seed search with algorithmic refinement to establish the performance upper bound. Multi-seed search generates diverse initial implementations to expand the exploration space. Algorithmic refinement then applies structural optimizations (e.g., operator fusion, algorithm replacement) to eliminate redundant computation and memory traffic at the source.
Subsequently, the hardware-specific tuning stage systematically realizes this performance potential through three sequential sub-stages with well-defined objectives: parallel mapping determines workload distribution across GPU cores, tensor tiling optimizes on-chip data reuse, and memory optimization refines data access patterns and pipeline execution. By decomposing the complex optimization into constrained sub-problems with explicit objectives, this staged approach effectively guides the LLM toward high-performance solutions.

To effectively support this workflow, we introduce a stage-aware multi-agent collaborative framework consisting of a code agent, a profile agent, and a debug agent. Within each stage, the three agents collaborate on focused objectives through structured information exchange (e.g., bottleneck analysis, fix suggestions). Across stages, context propagation retains only finalized optimization decisions and filters out intermediate exploration details. This design prevents objective interference across stages and ensures cumulative improvements throughout the optimization process.

The key contributions of this paper are as follows:
\begin{itemize}
\item We propose an expert-guided agent framework for GPU kernel generation that decomposes the optimization process into high-level algorithmic structure design and fine-grained hardware-specific tuning. By decomposing complex joint optimization into constrained sub-problems with explicit objectives, our approach effectively guides LLMs toward high-performance solutions.

\item We design a stage-aware multi-agent collaborative optimization mechanism, in which code, profile, and debug agents cooperate within each stage and propagate optimization decision across stages, achieving stable and accumulative improvements.

\item We systematically evaluate our approach on KernelBench and real-world workloads. Experimental results show that our approach achieves 100\% correctness while delivering up to $2.13\times$ speedup over PyTorch implementations, demonstrating the effectiveness of integrating expert knowledge into the kernel generation agent.
\end{itemize}

\section{Background and Related Work}
\subsection{GPU Kernel Optimization}
GPU kernel optimization relies on a set of widely applicable techniques: (1) \textit{operator fusion} reduces kernel launch overhead and redundant memory accesses by merging multiple operations; (2) \textit{parallel mapping} determines how computation is distributed across streaming multiprocessors (SMs), directly affecting hardware utilization; (3) \textit{tiling} partitions computation into blocks that fit register and cache capacities to maximize data reuse; (4) \textit{pipelining} overlaps computation with data movement to hide memory access latency.

While these optimization principles are largely shared across GPU platforms, their concrete implementations are highly hardware-dependent: SM configuration, register file size, on-chip memory hierarchy, and memory bandwidth all critically influence optimal design choices. High-performance kernels today predominantly come from expert-optimized libraries (e.g., cuBLAS, cuDNN) or specialized implementations (e.g., FlashAttention~\cite{flashattention}), achieving near-peak performance but at high development cost and limited generality across workloads and platforms.

\subsection{Compilers and Triton}
Compilers and domain-specific languages (DSLs) partially alleviate these challenges by providing higher-level abstractions while delegating low-level details to the compiler. Triton~\cite{triton}, a Python-embedded DSL with tile-based programming abstraction, has been widely adopted in modern LLM systems (e.g., PyTorch 2.0, vLLM) for implementing performance-critical kernels. However, achieving high performance with Triton still requires careful manual tuning of parallelism, tiling, and pipelining decisions—suboptimal configurations can lead to substantial performance loss. While CUDA offers finer-grained control, its larger optimization space makes automated optimization harder to constrain. Accordingly, this work focuses on Triton, as its structured abstraction provides a more tractable space for LLM-driven optimization.

\subsection{LLM-based Kernel Generation}
\textbf{Benchmarks and Datasets} Kernel generation benchmarks require both functional correctness and high performance, distinguishing them from general-purpose code generation benchmarks like HumanEval and MBPP. KernelBench provides a hierarchical benchmark consisting of 250 PyTorch workloads and proposes the $\mathrm{Fast}_p$ metric to jointly evaluate correctness and speedup. TritonBench presents a comprehensive Triton-specific benchmark comprising real-world Triton kernels collected from GitHub and operators with PyTorch reference implementations.

\textbf{Reinforcement Learning–based Methods} Pretrained LLMs struggle to generate high-performance GPU kernels due to limited understanding of hardware characteristics and optimization principles. To address this, recent work applies task-specific fine-tuning and reinforcement learning (RL). Kevin-32B improves kernel correctness and performance through multi-round RL training. AutoTriton and TritonRL apply supervised fine-tuning to learn Triton syntax, then use RL to further refine performance. However, these approaches require large-scale domain-specific datasets and substantial computational resources, limiting their practical applicability.

\textbf{Agent-based Methods} Multi-agent systems offer an alternative approach that avoids expensive training costs. CudaForge employs a hardware-feedback-driven framework composed of a Coder and a Judger, using Nsight Compute (NCU) metrics to guide optimization. AI CUDA Engineer uses stochastic mutation and iterative search to explore kernel designs. 
However, these methods lack holistic understanding of optimization principles and rely on random adjustments guided by coarse-grained performance feedback, leading to unstable behavior and limited improvements.

\begin{figure*}[tbhp]
    \centering
    \resizebox{1\linewidth}{!}{
     \includegraphics[width=1\linewidth]{figures/framework.pdf}
     }
     \caption{\textbf{Overview of XXXX.}
     }
     \label{fig:pipeline}
\end{figure*}

\section{Method}

\subsection{Overview}
In this section, we introduce XX, an expert-guided multi-agent framework for automatic generation of high-performance GPU kernels. As illustrated in Figure~\ref{fig:pipeline}, XX comprises two core components:
(1) an expert-guided staged optimization strategy that decomposes the kernel optimization space into well-defined sub-problems;
(2) a stage-aware multi-agent collaboration mechanism that coordinates specialized agents within each stage and propagates refined decisions across stages.
In the following parts, we first describe how we construct our expert-guided staged optimization. Then we detail how to design multi-agent system for each decomposed stage to progressively refine the GPU kernels.
% We first describe the construction of our expert-guided staged optimization strategy. Then, we detail the multi-agent system design that progressively refines GPU kernels through coordinated optimization at each stage.
% The following subsections describe each component of the framework in detail.

\subsection{Expert-Guided Staged Optimization}
The GPU kernel optimization space is highly complex, with optimization decisions at different levels tightly coupled. As a result, directly applying LLMs to perform end-to-end generation often struggles to achieve high performance.
To address this challenge, we introduce expert knowledge to decompose GPU kernel optimization into two levels: high-level algorithmic structure design and fine-grained hardware-specific tuning.
The algorithmic structure design determines the computational structure and dataflow organization, fundamentally establishing the upper bound of achievable performance.
The hardware-specific tuning governs the execution behavior and resource utilization on the target device, influencing the runtime performance.
This hierarchical decomposition reduces the decision complexity and improves the controllability of the optimization process.

\subsubsection{Algorithmic Structure Design}
Algorithmic structure design establishes a high-quality algorithmic foundation that determines the performance upper bound. We address a fundamental challenge: the optimal algorithmic approach for a given operator is often non-obvious, and LLM-generated initial implementations typically contain structural inefficiencies. We combine two complementary strategies: multi-seed search explores different algorithmic paradigms to provide implementation diversity, while algorithmic refinement optimizes each paradigm to eliminate inefficiencies.

\textbf{Multi-Seed Search}
For a given operator, fundamentally different algorithmic paradigms may exist. For example, convolution can be implemented as either direct nested-loop computation or im2col-based matrix multiplication. Since our structured optimization workflow constrains the decision space at each subsequent stage, the initial paradigm choice is critical to final performance. Therefore, we generate a small set of initial seeds with distinct algorithmic structures to provide diverse starting points. After applying algorithmic refinement to each seed, a Filter evaluates their performance and retains only the best-performing candidate for subsequent stages.

\textbf{Algorithmic Refinement}
While multi-seed search provides diverse algorithmic directions, each seed's initial implementation typically contains structural inefficiencies. Algorithmic refinement applies expert-guided transformations to optimize the computational structure within each seed's paradigm. Given a seed kernel, the LLM analyzes its operator composition and dataflow structure, then performs semantics-preserving optimizations such as operator fusion, algorithm replacement, or dataflow reorganization. These transformations reduce unnecessary memory accesses and computation at the algorithmic level.

As shown in Figure.~\ref{}, we take Attention computation as a representative example to illustrate the necessity of algorithmic refinement. Traditional implementations decompose Attention into multiple independent kernels: computing $QK^{T}$, performing softmax normalization, and multiplying with $V$. For a sequence of length $N$, this approach materializes the $N \times N$ attention matrix in global memory, incurring $O(N^2)$ additional memory traffic. FlashAttention restructures this by introducing an online softmax mechanism: during the traversal of $K$ and $V$, the kernel maintains normalization statistics online, completing attention score computation, normalization, and weighted accumulation within a single fused kernel, thereby eliminating this memory bottleneck.

Overall, this combination balances coarse-grained paradigm exploration with fine-grained structural refinement, establishing a strong algorithmic foundation for subsequent hardware-specific tuning.

\subsubsection{Hardware-Specific Tuning}
After establishing the high-level algorithmic structure, we perform fine-grained hardware-specific tuning through three progressive stages: (1) \textit{parallel mapping} determines parallel strategy to fully utilize GPU cores; (2) \textit{tensor tiling} selects block-level tiling granularity to maximize on-chip data reuse; and (3) \textit{memory optimization} optimizes data access patterns and execution pipelines. 
Each stage constrains the decision space for subsequent stages. Parallel mapping fixes the global parallel structure; tiling then operates within individual program instances under this fixed structure; memory optimizations further refine execution under both fixed parallelization and tiling. By decomposing the complex joint optimization into a sequence of constrained sub-problems, we provide the LLM with well-scoped objectives at each stage, enabling more effective and focused performance improvement.
% After establishing the high-level algorithmic design, the workflow proceeds to the fine-grained hardware-specific tuning stage. 
% This stage applies progressive optimizations tailored to the GPU architecture's parallelism and memory hierarchy through three sub-stages: (1) determining parallel mapping strategy to fully exploit hardware parallelism; (2) selecting appropriate block-level tiling granularity to improve on-chip data reuse; and (3) optimizing data access patterns and pipeline execution under fixed parallel structure and tiling choices. By explicitly modeling this progressive workflow, we provide the LLM with a clear and constrained optimization trajectory, enabling it to focus on well-defined optimization objectives at each stage and progressively approach achievable performance limits.

\textbf{Parallel Mapping} The parallel mapping stage determines how computational tasks are distributed across GPU streaming multiprocessors (SMs). The key optimization objective is to identify parallelizable dimensions of the operator (e.g., batch, head, expert) and map them to the GPU execution grid to fully utilize multi-core parallelism.

Formally, we define parallel mapping as a dimension-to-grid mapping:
  \begin{equation}
  (d_1, d_2, \ldots, d_n) \rightarrow (G_0, G_1, G_2)
  \end{equation}
where $d_i$ denotes a parallelizable dimension of the operator, and $G_i$ is the number of program instances launched in the $i$-th grid dimension. For example, multi-head attention with batch size $B$ and $H$ heads can be mapped as $(B, H) \rightarrow (G_0, G_1) = (B, H)$, launching $B \times H$ parallel instances.
% By explicitly modeling the mapping between operator parallel dimensions and the execution grid, we abstract parallel structure design as high-level decisions with clear semantic constraints. This mapping decision fixes the parallel structure for subsequent stages: all further optimizations (tiling, memory access) operate within individual program instances under this established parallel decomposition.
  
\textbf{Tensor Tiling}
After determining the global parallel structure, the tensor tiling stage focuses on designing computation granularity within each program instance. The core objective is to determine the size and shape of tensor tiles processed by a single program instance in one execution iteration.

We formalize tensor tiling as a dimension-to-block mapping:
\begin{equation}
  (d_1, d_2, \ldots, d_m) \rightarrow (B_1, B_2, \ldots, B_m)
  \end{equation}
where $d_i$ denotes the size of operator dimension $i$, and $B_i$ represents the tile size for that dimension. The relationship between parallel mapping and tiling is given by $G_i = \lceil d_i / B_i \rceil$: for a dimension of size $d_i$, selecting tile size $B_i$ determines the required number of program instances $G_i$ to cover the entire dimension.

This stage directly determines how data resides in registers and shared memory, thereby governing the efficiency trade-off between computation and memory access. Larger tiles improve on-chip data reuse and reduce global memory access frequency, but increase register and shared memory usage, potentially limiting parallelism. Conversely, smaller tiles have lower resource overhead but cannot fully exploit data locality. 
Based on this trade-off, tile size selection can be modeled as a constrained optimization problem:
  \begin{equation}
  \begin{aligned}
  &(B_1^*, \ldots, B_m^*) = \arg\max_{(B_i) \in \mathcal{C}} \Phi(B_1, \ldots, B_m) \\&\text{s.t.} \quad \text{RegisterUsage}(B_i) \leq \text{MaxRegisters} \\
  &\phantom{\text{s.t.}} \quad \text{SharedMemUsage}(B_i) \leq \text{MaxSharedMem}
  \end{aligned}
  \end{equation}
where $\Phi(\cdot)$ measures computational efficiency and the constraints ensure resource usage does not exceed hardware limits.

\textbf{Memory Optimization}
The memory optimization stage operates under fixed parallel structure and tiling configuration from prior stages. The stage-specific prompt directs the agent to optimize two complementary aspects: (1) memory access patterns, organizing global memory accesses into coalesced loads to improve bandwidth utilization; and (2) software pipelining, adjusting multi-buffering depth to overlap data movement with computation and hide memory latency. These optimizations target memory system bottlenecks while preserving the established computational structure, enabling performance gains without architectural disruption.

\subsection{Stage-Aware Multi-Agent Collaboration}

After decoupling the optimization process into distinct stages, we introduce a stage-aware multi-agent collaboration mechanism to perform optimization within each stage. A single agent responsible for code generation, performance analysis, and debugging across all stages often exhibits unstable optimization behavior, as accumulated cross-stage context obscures the current objective and may lead to repeated revisions or regressions of earlier optimizations. To mitigate these issues, our framework adopts a structured multi-agent design with explicit context flow management, decomposing responsibilities across specialized agents while enabling stable and accumulative optimization.

\subsubsection{Multi-Agent Design}
The multi-agent system decomposes the optimization process into functionally distinct roles that collectively support stable and effective optimization. At each stage, three specialized agents form a closed-loop workflow around a unified optimization objective:
\begin{itemize}
\item \textbf{Profile Agent:} Takes profiling metrics (e.g., NVIDIA Nsight Compute (NCU) reports), the current kernel code, and the stage objective as input, and identifies the primary performance bottleneck at the current stage (e.g., compute-bound, memory-bound, or parallelism-limited).

\item \textbf{Code Agent:} Takes optimization suggestions and the current kernel implementation as input, and generates or locally modifies kernel code to address the current stage's optimization objective.

\item \textbf{Debug Agent:} Takes error logs and the current kernel code as input, diagnoses compilation errors, runtime exceptions, or numerical inconsistencies, and provides targeted fixes.
\end{itemize}

\subsubsection{Context Flow Management}
Context flow management governs how optimization context is organized, filtered, and shared to ensure effective collaboration among agents and coherent optimization progress. We introduce two complementary mechanisms: inter-stage context propagation, which manages context across stage transitions to support cumulative optimization, and intra-stage information exchange, which coordinates agent interactions within each stage for focused exploration.

\textbf{Inter-Stage Context Propagation.}
Inter-stage context propagation keeps agents focused on the current stage objective while preserving finalized decisions from previous stages. At stage boundaries, we filter and reorganize accumulated context to avoid cross-stage interference. When transitioning from stage $t$ to $t{+}1$, the system: (1) retains only finalized decisions and kernel implementations from stage $t$, discarding intermediate exploration artifacts; (2) constructs a new context view that explicitly defines the optimization target for stage $t{+}1$ (e.g., ``current stage: grid\_and\_parallel, focus on grid layout and parallelism optimization''); and (3) propagates this context to all agents so that the Code Agent builds incrementally upon prior work while the Profile Agent focuses on metrics relevant to the new objective. This mechanism establishes a stable, cumulative optimization trajectory without regressing earlier improvements.

\textbf{Intra-Stage Information Exchange.}
Intra-stage information exchange coordinates efficient collaboration among agents through structured JSON interfaces that compress and isolate context. The profile agent outputs bottleneck analysis, optimization method, and modification plan. When failures occur, the debug agent outputs critical issue identification and required fixes. The code agent consumes these structured inputs along with current kernel code and stage-specific constraints to generate modifications. Figure~\ref{fig:agent_exchange} illustrates this information flow with concrete examples. This design avoids context overload and enables tight collaboration loops for efficient iterative optimization.

Overall, combining multi-agent collaboration with structured context flow management enables progressive, cumulative optimization across stages while maintaining efficient, focused exploration within each stage.


\section{Experiments}
In this section, we comprehensively evaluate EGG through systematic experiments on Triton kernel generation to analyze its effectiveness and behavior.

\subsection{Experimental Setup}
\textbf{Hardware Platforms} We evaluate EGG on two representative GPU platforms to validate its generality:
(1) NVIDIA GeForce RTX 4090 with 24GB GDDR6X memory, Ada Lovelace architecture, 128 SMs and 16,384 CUDA cores;
(2) NVIDIA GeForce RTX 5090 with 32GB GDDR6X memory, Blackwell architecture, 170 SMs and 21,760 CUDA cores;
For brevity, we present results on RTX 4090, 
% in the main text, 
while leaving RTX 5090 results in Appendix, which also demonstrats consistent performance trends.


\textbf{Software} All experiments are conducted with CUDA 13.0 and Triton 3.5.1. LLM inference is primarily supported by GPT-5.1. \hyc{（Results using Claude Sonnet 4.5 are provided in the Appendix.）} Multi-seed search uses two seeds to balance optimization quality and computational cost.

\textbf{Benchmarks}
We adopt KernelBench\cite{ouyang2025kernelbench} as the primary evaluation benchmark. It contains 250 tasks across three difficulty levels: Level 1 includes 100 basic operators (e.g., matrix multiplication, ReLU); Level 2 contains 100 medium-difficulty tasks combining multiple operators (e.g., fused GEMM+ReLU+Add); Level 3 comprises 50 challenging tasks covering complete neural network architectures (e.g., AlexNet, ResNet). Each task provides a reference PyTorch implementation along with standardized input and output specifications for reliable correctness and performance evaluation.

\textbf{Baselines}
To evaluate the performance of EGG, We compare against the following baselines:
(1) ChatGPT-5.1\cite{openai_chatgpt_5_1}, a state-of-the-art commercial LLM with advanced reasoning capabilities;
(2) DeepSeek-V3.2\cite{liu2025deepseek}, a high-performance open-source model;
(3) PyTorch Eager, the default execution mode that invokes vendor-optimized libraries (cuBLAS, cuDNN) with operator-by-operator execution;
(4) Torch Compile\cite{ansel2024pytorch}, PyTorch's graph compilation framework that applies operator fusion, memory planning, and optimized kernel selection;
(5) AutoTriton\cite{li2025autotriton}, an LLM fine-tuned with reinforcement learning specifically for Triton programming tasks;
(6) CudaForge\cite{zhang2025cudaforge}, a multi-agent system that performs performance-feedback-driven iterative refinement for kernel generation.

\textbf{Metrics}
\hyc{Following previous works~\cite{ouyang2025kernelbench,li2025tritonbench,baronio2025kevin}, we evaluate all methods using the following metrics:}
(1) Success Rate: the fraction of tasks that successfully compile and pass correctness verification;
(2) $\mathrm{Fast}_1$ Rate: the fraction of tasks for which the generated kernels outperform PyTorch Eager (speedup $>1.0\times$);
(3) Speedup: the average execution time improvement of generated kernels over PyTorch Eager, computed only for successful kernels.

\subsection{Overall Performance}

Table~\ref{tab:kernelbench_overall} presents comprehensive performance comparison across KernelBench's three difficulty levels. EGG achieves 100\% success rate across all levels while delivering substantial performance improvements: 2.13$\times$ speedup over PyTorch Eager (averaged across all levels) and 1.60$\times$ over the auto-compilation baseline Torch Compile.

\textbf{Performance across difficulty levels} EGG consistently outperforms all baselines across all difficulty levels. In Level 1 basic operators,  we achieve 1.83$\times$ speedup with 72\% $\mathrm{Fast}_1$ rate and 100\% success rate. As task complexity increases to Level 2 fused operators, our advantage becomes more pronounced: we achieve 2.73$\times$ speedup with 100\% $\mathrm{Fast}_1$ rate, where every generated kernel exceeds PyTorch Eager performance, demonstrating the effectiveness of our expert-guided staged optimization. In the most challenging Level 3 end-to-end models, our method maintains 100\% success rate while achieving 1.52$\times$ speedup and 94\% $\mathrm{Fast}_1$ rate, showing robust performance even in the most complex tasks.

\textbf{Comparison with baselines} General-purpose LLMs (DeepSeek, ChatGPT) and the domain-specialized AutoTriton achieve only 0.30-0.53$\times$ speedup with 36-60\% success rates, indicating that neither general reasoning capabilities nor domain-specific fine-tuning alone suffices for systematic kernel optimization.

Torch Compile, leveraging graph-level optimization and pre-built kernel libraries, achieves 1.09-1.38$\times$ speedup. However, its reliance on existing kernel implementations limits performance on novel operator compositions. \hyc{In contrast, EGG synthesizes custom kernels tailored to specific workload patterns}, delivering 1.68$\times$, 1.98$\times$, and 1.12$\times$ additional improvements across the three levels, with the largest gain on Level 2 fused operators.

CudaForge employs multi-agent iterative refinement with hardware feedback, achieving 1.28-2.10$\times$ speedup. However, without explicit guidance on optimization direction, it relies on trial-and-error exploration, leading to unstable results on complex tasks (e.g., Level 3: 68\% $\mathrm{Fast}_1$ rate). Our approach provides clear optimization focus for each stage (e.g., parallelization, memory access), enabling more consistent improvements and superior final performance.


\textbf{Cost efficiency} EGG completes generation for a single task in approximately 20 minutes on a single RTX 4090 GPU, consuming around 50,000 tokens per kernel. This demonstrates that decomposing optimization into explicit stages enables high-quality implementations at reasonable computational cost.

\subsection{Ablation Study}
  \begin{table}
  \centering
  \caption{Ablation study on multi-seed search, algorithmic refinement and hardware-specific tuning.}
  \label{tab:ablaltion}
\begin{tabular}{ccccc}
\toprule
\begin{tabular}[c]{@{}c@{}}\textbf{Multi-}\\ \textbf{Seed}   \end{tabular} & \begin{tabular}[c]{@{}c@{}}\textbf{Algo.}\\ \textbf{Refinement}\end{tabular} & \begin{tabular}[c]{@{}c@{}}\textbf{HW}\\ \textbf{Tuning}\end{tabular} & $\mathbf{Fast}_1$ & \textbf{Speedup} \\
\midrule
  \ding{51}   &       \ding{51}       &        \ding{51}       &   87.6\%   &  2.13\times      \\
            & \ding{51}       &     \ding{51}       &   78.4\%   &   1.84\times     \\
  &        \ding{51}               &                                     &  60.8\%   &   1.52\times  \\
     &                      &        \ding{51}             &  72\%   &   1.46\times  
\end{tabular}
\bottomrule
\end{table}

\begin{figure}[t]
	\centering
	\includegraphics[width=3.3in]{figures/result.pdf}
	\caption{Cumulative speedup across four optimization stages.}
	\label{cul}
\end{figure}

As shown in Table~\ref{tab:ablaltion}, we systematically analyze the contribution of each component in our framework through ablation studies. 

\textbf{Multi-seed search} 
Introducing 2-seed search improves average speedup from 1.84$\times$ to 2.13$\times$ and $\mathrm{Fast}_1$ rate from 78.4\% to 87.6\%. Notably, the improvement varies across difficulty levels: Level 2 and Level 3 achieve 1.7$\times$ and 2.0$\times$ gains respectively, while Level 1 shows only 1.09$\times$ improvement.
This is because complex tasks (fused operators, full models) have larger algorithmic design spaces where diverse seeds enable broader exploration, while simple operators yield similar structures across different seeds. \hyc{This validates that multi-seed search effectively prevents early convergence, with benefits increasing with task complexity.}

\textbf{Stage-specific contributions} \hyc{Enabling only the algorithmic refinement stage achieves $1.52\times$ speedup but only 60.8\% $\mathrm{Fast}_1$ rate, indicating that while algorithmic refinement unlocks high performance potential, it fails to deliver consistent speedups across different kernels without hardware-aware constraints. In contrast, enabling only hardware-specific tuning yields 1.46$\times$ speedup with 72\% $\mathrm{Fast}_1$ rate, achieving more consistent acceleration across tasks but limited peak performance due to suboptimal initial algorithms.} Combining both components achieves the optimal performance, demonstrating their complementary nature: algorithmic refinement establishes the performance ceiling, while hardware-specific tuning ensures stable realization of that potential.


\textbf{Cumulative stage effects} Fig.~\ref{cul} illustrates the cumulative performance improvement across optimization stages. Algorithmic refinement stage (Stage 1) achieves 2.0$\times$ speedup over the initial seed, effectively establishing the performance baseline. Subsequent hardware-specific tuning stages (Stage 2-4) contribute an additional $1.7\times$ improvement through parallel mapping (1.33$\times$), tensor tiling (1.08$\times$), and memory optimization (1.13$\times$). While individual stage gains are modest, their cumulative effect remains consistent across operators, validating our expert-guided multi-stage optimization workflow.

\subsection{Case Study: 3D Transposed Convolution}

\begin{table}
\centering
\caption{Optimization trajectory for grouped 3D transposed convolution: latency and speedup relative to PyTorch eager mode.}
\label{case}
\begin{tabular}{lcc}
\toprule
\textbf{Stage} & \textbf{Latency (ms)} & \textbf{Speedup} \\
\midrule
PyTorch Eager                & 20.78 & --    \\
\midrule
Seed 1                       & 29.62 & 0.70$\times$  \\
Seed 2                       & 76.04 & 0.27$\times$  \\
\midrule
Algorithmic Refinement (S1)  & 19.41 & 1.07$\times$  \\
Algorithmic Refinement (S2)  & 12.38 & 1.68$\times$  \\
\midrule
Parallel Mapping                 & 10.78 & 1.93$\times$   \\
Tensor Tiling                 &  9.20 & 2.26$\times$ \\
\bottomrule
\end{tabular}
\end{table}

We demonstrate EGG's complete optimization workflow through a representative example: grouped 3D transposed convolution with batch size = 16, input/output channels = 32, kernel size = (3, 5, 7), stride = (2, 2, 2), padding = (1, 2, 3), and groups = 4. On an RTX 4090 GPU, the PyTorch Eager baseline achieves 20.78 ms.

\textbf{Multi-seed search} The system generates two candidate implementations with different algorithmic structures. Seed 1 adopts an output-centric backward-mapping strategy, computing corresponding inputs for each output position through reverse indexing, achieving 29.62 ms (0.7$\times$ speedup vs PyTorch). Seed 2 uses a direct nested iteration over input and kernel elements, resulting in 76.04 ms (0.27$\times$) due to poor data reuse. Both initial implementations underperform PyTorch baseline.

\textbf{Algorithmic refinement stage} The profile agent analyzes NCU metrics and identifies critical bottleneck. For Seed 1, backward-mapping requires enumerating all $C_{In} \times K_D \times K_H \times K_W$ combinations using expensive division/modulo operations in inner loops. When stride $>1$, this generates numerous invalid mappings, resulting in excessive control-flow overhead and warp divergence. The profile agent recommends restructuring to a forward-mapping scheme, where input elements drive iteration and output coordinates are computed via simple multiply-add with in-block accumulation. After the coder agent's modification, latency improves to 19.41 ms (1.07$\times$), surpassing PyTorch.

For Seed 2, the profile agent identifies extremely low arithmetic intensity and insufficient data reuse. It recommends transforming transposed convolution into matrix multiplication via im2col. This algorithmic restructuring dramatically reduces latency to 12.38 ms (1.68$\times$). Given its superior scalability, subsequent optimization focuses exclusively on this kernel.

\textbf{Hardware-Specific tuning stage}
\emph{Parallel Mapping.} NCU profiling reveals low SM utilization. The profile agent guides the code agent to remap output pixels, output channels, and groups across three \texttt{program\_id} axes, improving workload distribution. Latency reduces to 10.78 ms (1.93$\times$).
\emph{Tensor Tilling.}
Analyzing register and shared memory usage, the profile agent directs autotuning over \texttt{BLOCKM/N/K} configurations. After searching $(32,32,32)$, $(64,32,32)$, and $(32,64,32)$, the optimal configuration achieves 9.20 ms (2.26$\times$).

Table~\ref{case} summarizes the complete optimization trajectory. Starting from the initial 0.27$\times$ implementation (Seed 2), algorithmic refinement delivers 6.2$\times$ improvement to 1.68$\times$, establishing a solid computational foundation. Subsequent hardware-specific stages contribute cumulative 1.35$\times$ gain through parallel mapping (1.15$\times$) and tensor tiling (1.17$\times$). This demonstrates our framework's effectiveness: algorithmic refinement establishes performance potential, while staged hardware tuning systematically realizes it—achieving final 2.26$\times$ speedup over PyTorch and 8.4$\times$ over the initial implementation.



\subsection{Practical Application Verification}
\begin{table}
\centering
\caption{Speedup over PyTorch for real-world kernels from TritonBench.}
\label{tab:real_world_speedup}
\begin{tabular}{lcc}
\toprule
\textbf{Task} & \textbf{TritonBench} & \textbf{Ours} \\
\midrule
Flash Attention        & 1.16 & 1.90 \\
RoPE Embedding         & 1.86 & 3.56 \\
INT8 Dequant MatMul    & 1.23 & 1.73 \\
\bottomrule
\end{tabular}
\end{table}
To evaluate practical deployment effectiveness, we assess EGG on representative operators from TritonBench\cite{li2025tritonbench} that are commonly used in real-world LLM inference workloads, as shown in Table~\ref{tab:real_world_speedup}. These operators are derived from production Triton kernels in public GitHub repositories. Since the original benchmarks provide natural language descriptions as input, we convert them to PyTorch reference implementations to align with practical kernel optimization scenarios.

We compare our generated kernels against the original hand-written Triton implementations. EGG achieves substantial speedups across all operators: 1.64$\times$ for Flash Attention, 1.91$\times$ for RoPE Embedding, and 1.41$\times$ for INT8 Dequant MatMul. \hyc{These results demonstrate that expert manual tuning does not always achieve optimal performance for complex operators. EGG effectively harnesses LLMs' exploration capabilities to discover implementations that surpass hand-tuned production kernels, validating the practical value of our approach.}
% These results demonstrate that our expert-guided staged optimization can discover kernel implementations that outperform even hand-tuned kernels from production systems.

\section{Conclusion}
In this work, we propose an expert-guided GPU kernel generation agent framework that structures optimization into algorithmic structure design and hardware-specific tuning, systematically guiding LLM decisions with domain-specific knowledge. A stage-aware multi-agent collaboration mechanism enables code, profile, and debug agents to achieve stable, cumulative improvements. Experiments on KernelBench and real Triton workloads demonstrate significant performance gains while maintaining 100\% correctness, outperforming existing agent-based and reinforcement learning methods.
% In this work, we propose an expert-guided GPU kernel generation agent framework. The framework structures optimization into two levels: algorithmic restructuring and hardware-specific optimization, systematically guiding LLM decisions with expert knowledge. At the algorithmic level, constrained multi-seed search explores computational structures and identifies high-performance patterns. At the hardware level, it sequentially optimizes parallel mapping, tensor tiling, and memory access. A phase-aware multi-agent collaboration mechanism enables coding, performance analysis, and debugging agents to work cooperatively, achieving stable and cumulative improvements. Experiments on KernelBench and real Triton workloads show that our framework generates significantly faster kernels while maintaining 100\% correctness, outperforming existing agent-based and reinforcement learning methods.
 
\nocite{langley00}

\bibliography{example_paper}
\bibliographystyle{icml2026}

%%%%%%%%%%%%%%%%%%%%%%%%%%%%%%%%%%%%%%%%%%%%%%%%%%%%%%%%%%%%%%%%%%%%%%%%%%%%%%%
%%%%%%%%%%%%%%%%%%%%%%%%%%%%%%%%%%%%%%%%%%%%%%%%%%%%%%%%%%%%%%%%%%%%%%%%%%%%%%%
% APPENDIX
%%%%%%%%%%%%%%%%%%%%%%%%%%%%%%%%%%%%%%%%%%%%%%%%%%%%%%%%%%%%%%%%%%%%%%%%%%%%%%%
%%%%%%%%%%%%%%%%%%%%%%%%%%%%%%%%%%%%%%%%%%%%%%%%%%%%%%%%%%%%%%%%%%%%%%%%%%%%%%%
\newpage
\appendix
\onecolumn
\section{GPU Kernel Optimization Architect Prompt}

This section presents the complete prompt used to guide the LLM in identifying high-level algorithmic optimizations for GPU kernels.

\begin{mdframed}[linewidth=1pt, linecolor=black, backgroundcolor=gray!5]
\begin{lstlisting}
You are a GPU kernel optimization architect. Analyze the kernel and identify
**ONE high-level algorithmic optimization**.

# PyTorch Reference
```python
$python_code
```

# Current Triton Kernel
```python
$cuda_code
```

$performance_section

---

## Analysis Steps

1. **Code Analysis**: Count kernels, identify operations, check for inefficiencies
2. **Performance Diagnosis**: Use metrics/latency to identify bottleneck type
3. **Root Cause**: Combine code + performance to find the core issue

## Optimization Categories (pick ONE if worth optimizing):

### 1. Operator Fusion
Fuse consecutive ops into fewer kernels to reduce memory traffic and launch overhead.

### 2. Algorithm Replacement
Replace naive algorithm with optimized variant.
- For Attention: Flash Attention, online softmax
- For Convolution: Winograd, im2col
- **For RNN/GRU/LSTM**: Persistent kernel with HYBRID computation
  - **CRITICAL**: Use hybrid approach for best performance:
    * Precompute input-side gates ONCE (outside kernel): `gates_x = (T*B, In) @ W_ih`
    * Persistent kernel (inside): only recurrent-side: `for t: gates_h = h @ W_hh`
  - Time loop `for t in range(T)` must be inside kernel, NOT in Python
  - Launch kernel once per layer, not once per timestep
  - Expected speedup: 10-100x (vs per-timestep launches)

### 3. Kernel Launch Reduction
Combine multiple small kernels to reduce overhead.
- **For RNN/GRU/LSTM**: See "Algorithm Replacement" above for persistent kernel approach

### 4. Memory Layout Optimization
Use in-place operations, buffer reuse, or better layouts.

## Should We Optimize?

Before proposing optimization, determine if it's worthwhile:
- **Not worth optimizing** if:
  - Code is already near-optimal (expected speedup < 10%)
  - Bottleneck cannot be addressed (hardware limited, already optimal algorithm)
  - Optimization would add significant complexity with minimal gain

- **Worth optimizing** if:
  - Clear algorithmic inefficiency exists (multiple kernels, suboptimal algorithm)
  - Expected speedup >= 20%
  - Concrete optimization path available

## Output (JSON)

```json
{
  "worth_optimizing": "yes/no",
  "reason": "<Why worth or not worth optimizing, 1 sentence>",
  "bottleneck": "<Root cause in 1-2 sentences, empty if not worth optimizing>",
  "optimisation method": "<Specific optimization in 1-2 sentences, empty if not worth optimizing>",
  "modification plan": "<Implementation steps in 2-3 sentences, empty if not worth optimizing>",
  "expected_speedup": "<e.g., '30-40%', empty if not worth optimizing>"
}
```

Return JSON only.
\end{lstlisting}
\end{mdframed}
%%%%%%%%%%%%%%%%%%%%%%%%%%%%%%%%%%%%%%%%%%%%%%%%%%%%%%%%%%%%%%%%%%%%%%%%%%%%%%%
%%%%%%%%%%%%%%%%%%%%%%%%%%%%%%%%%%%%%%%%%%%%%%%%%%%%%%%%%%%%%%%%%%%%%%%%%%%%%%%

\end{document}

% This document was modified from the file originally made available by
% Pat Langley and Andrea Danyluk for ICML-2K. This version was created
% by Iain Murray in 2018, and modified by Alexandre Bouchard in
% 2019 and 2021 and by Csaba Szepesvari, Gang Niu and Sivan Sabato in 2022.
% Modified again in 2023 and 2024 by Sivan Sabato and Jonathan Scarlett.
% Previous contributors include Dan Roy, Lise Getoor and Tobias
% Scheffer, which was slightly modified from the 2010 version by
% Thorsten Joachims & Johannes Fuernkranz, slightly modified from the
% 2009 version by Kiri Wagstaff and Sam Roweis's 2008 version, which is
% slightly modified from Prasad Tadepalli's 2007 version which is a
% lightly changed version of the previous year's version by Andrew
% Moore, which was in turn edited from those of Kristian Kersting and
% Codrina Lauth. Alex Smola contributed to the algorithmic style files.
